Der Kühlschrank pumpt Wärme vom Kühlraum ($T_\text{k} = \Tk$) in die Küche ($T_\text{w} = \Tw$). Die maximal mögliche Leistungsziffer ist (FoTaBe S.~172)
\begin{align}
 \eta_\text{K,max} &= \frac{T_\text{k}}{T_\text{w}-T_\text{k}}
\end{align}
der Kühlschrank erreicht davon den Anteil $f = \frO$: $\eta_\text{K} = f\,\eta_\text{K,max} = \eK$. Mit $\eta_\text{K} = Q_\text{k}/W$ und der elektrischen Energie $W = P\,t$ ist die entzogene Wärme
\begin{align}
 Q_\text{k} &= \eta_\text{K}\,P\,t = \QkF \\
   &= \frO\cdot\frac{\Tk}{\Tw-\Tk}\cdot\Pk\cdot\td \\
   &= \Qk \approx \result{\QkP{6}{2}}
\end{align}
